\documentclass[11pt]{article}
\usepackage[a4paper,margin=1in]{geometry}
\usepackage[T1]{fontenc}
\usepackage{lmodern}
\usepackage{amsmath}
\setlength{\parindent}{0pt}
\setlength{\parskip}{8pt}

\begin{document}

\section*{Part C: Interpretation}

Reputation satisfaction ($X_1$) appears most strongly associated with
preferred provider. Among respondents who prefer B, 65.0\% report high
reputation satisfaction, compared with 39.0\% for A and 31.7\% for C.
The corresponding differences across providers are smaller for price
satisfaction ($X_2$) and service-variety satisfaction ($X_3$).
Thus, reputation provides the clearest individual distinction among
the three preference groups in this sample.

The model comparison also indicates that the attributes retain substantial
dependence after conditioning on preferred provider. At $\alpha=12$, the
fully dependent model M1 has the highest log marginal likelihood and
receives essentially all posterior model probability among the five
candidates under equal model priors. However, all five models achieve
the same in-sample accuracy of $450/877\approx51.31\%$ under our fixed
tie-breaking rule. Their probability estimates differ, but their predicted
classes are mostly unchanged. For the high/high/low profile, M1 ties
between A and B; the equal observed counts for these two classes mean
that consistently choosing either gives the same number of correct
classifications. Stronger model evidence therefore does not translate
into higher classification accuracy here.

For provider A, the reputation gap relative to B suggests a useful focus
for further customer research. A could investigate which aspects of its
reputation customers value and how reputation interacts with perceptions
of price and service variety. These results identify a potential area
for investigation, rather than demonstrating that a reputation campaign
would increase preference for A.

Provider C has the lowest proportions reporting high reputation and
service-variety satisfaction. Moreover, none of the five models predicts
C as the most probable class for any of the eight attribute profiles,
although 164 respondents prefer C. This suggests that the three binary
attributes do not adequately identify a segment in which C is the leading
choice. Collecting additional attributes or more detailed satisfaction
measures could help explain C's appeal and improve its identification.

These findings describe associations, not causal effects. The analysis
does not establish that changing an attribute would change provider
preference. In addition, accuracy was evaluated on the fitting data;
performance on independent observations remains to be assessed.

\end{document}
